\begin{solapartat}
\punts{0,1} $A = 0,02\ \mathrm{m}$

\smallskip
\punts{0,2} $T = \dfrac{0,04\ \mathrm{s}}{5} = 0,008\ \mathrm{s} \Rightarrow f = \dfrac{1}{T} = 125\ \mathrm{s^{-1}}$

\smallskip
\punts{0,1} $x(t) = A\cos(2\pi f t + \theta_0)$

\smallskip
\punts{0,2} $x(t) = 0,02\ \mathrm{m}\cos\left(250\pi\,\dfrac{\mathrm{rad}}{\mathrm{s}}\,t\right)$ \ o bé: \ $x(t) = 0,02\ \mathrm{m}\sin\left(250\pi\,\dfrac{\mathrm{rad}}{\mathrm{s}}\,t + \dfrac{\pi}{2}\ \mathrm{rad}\right)$

\smallskip
\punts{0,2}
\figura[0.5]{sol_fig1.png}

\punts{0,2} Ara, $f' = 62,5\ \mathrm{s^{-1}}$; \ $T' = \dfrac{1}{62,5} = 0,016\ \mathrm{s}$
\end{solapartat}

\begin{solapartat}
\begin{description}[leftmargin=4.5em, labelwidth=4em, labelsep=0.5em, itemsep=2pt]
  \item[\punts{0,3}b1)] Ni la freqüència ni la longitud d'ona canviaran.
  \item[\punts{0,3}b2)] La freqüència no canviarà, però la longitud d'ona augmentarà.
  \item[\punts{0,4}b3)] Percebrem una freqüència major, i la longitud d'ona (distància entre fronts successius) disminuirà.
\end{description}
\end{solapartat}
